\section{Experiments}
\subsection{Protocol}
LIBERO-GOAL and LIBERO-OBJECT \citep{liu2023libero} use official suite-specific
six-task pretrained initializations, followed by four continual tasks trained
for 100 epochs each. Seeds are 41, 42, and 43. All seen continual tasks are
evaluated after each stage on 20 initial states per task, yielding 200 formal
episodes per run. SHBR collects 40 additional training episodes per run.
Evaluation uses a 600-step horizon, five initial zero actions, and maximum-logit
component selection followed by Normal sampling. Latest $C/Q/S$ comparisons
use paired noise as bound in their evidence. Historical baselines do not
assert identical trajectory-level action noise.
\subsection{Metrics}
Let $p_{s,t}$ denote success on task $t$ after stage $s$, relabeling the four
continual tasks as $1,\ldots,T$. We use
\begin{align}
\mathrm{AUC}&=\frac1T\sum_{t=1}^T\frac1{T-t+1}\sum_{s=t}^T p_{s,t},\\
\mathrm{NBT}&=\frac1{T-1}\sum_{t=1}^{T-1}\frac1{T-t}
\sum_{s=t+1}^T(p_{t,t}-p_{s,t}).
\end{align}
FWT is mean diagonal success, not gain over a zero-shot baseline.
OldFinal averages final old-task success; FinalAll averages all final tasks.
NBT is lower-better and can be negative.
\subsection{Overall comparison}
\begin{table}[tb]
\centering\small
\caption{GOAL: reproduced baseline and complete method. Percent; AUC is mean $\pm$ sample standard deviation over seeds 41--43; other columns are means.}\label{tab:main-GOAL}
\begin{tabular}{lrrrr}
\toprule
Method & AUC $\uparrow$ & FWT $\uparrow$ & NBT $\downarrow$ & FinalAll $\uparrow$ \\
\midrule
MLR+IFA & $76.35\pm6.38$ & 76.67 & 0.19 & 76.67 \\
+ RA-Coreset & $81.98\pm3.00$ & 83.75 & 2.59 & 80.83 \\
+ RA + ALIGN & $81.11\pm3.42$ & 85.83 & 10.83 & 76.25 \\
Complete method & $87.26\pm2.65$ & 89.58 & 5.28 & 85.00 \\
\bottomrule
\end{tabular}
\end{table}

\begin{table}[tb]
\centering\small
\caption{OBJECT: reproduced baseline and complete method. Percent; AUC is mean $\pm$ sample standard deviation over seeds 41--43; other columns are means.}\label{tab:main-OBJECT}
\begin{tabular}{lrrrr}
\toprule
Method & AUC $\uparrow$ & FWT $\uparrow$ & NBT $\downarrow$ & FinalAll $\uparrow$ \\
\midrule
MLR+IFA & $68.30\pm3.12$ & 66.67 & -2.78 & 65.42 \\
+ RA-Coreset & $73.23\pm1.78$ & 73.33 & -0.09 & 69.58 \\
+ RA + ALIGN & $80.17\pm3.66$ & 77.50 & -4.91 & 75.83 \\
Complete method & $86.49\pm2.19$ & 80.42 & -11.02 & 85.00 \\
\bottomrule
\end{tabular}
\end{table}

The complete method improves mean AUC over reproduced MLR+IFA by 10.90
points on GOAL and 18.19 on OBJECT. These complete-method comparisons
are complemented by the controls below.
\draftnote{Complete the baseline protocol table; separate paper-reported
external results from reproduced results.}
\subsection{Retention and update controls}
\begin{table}[tb]
\centering\small
\caption{GOAL: update and retention controls. AUC shows mean $\pm$ sample standard deviation; final scores are three-seed means.}\label{tab:controls-GOAL}
\begin{tabular}{lrrr}
\toprule
Method & AUC $\uparrow$ & OldFinal $\uparrow$ & FinalAll $\uparrow$ \\
\midrule
Expert control $C$ & $82.50\pm7.10$ & 80.56 & 80.00 \\
Optimizer-matched $C$ & $83.72\pm3.35$ & 77.22 & 80.00 \\
Ungated $Q$ & $83.51\pm1.46$ & 78.33 & 80.42 \\
Complete method $S$ & $87.26\pm2.65$ & 83.89 & 85.00 \\
\bottomrule
\end{tabular}
\end{table}

\begin{table}[tb]
\centering\small
\caption{OBJECT: update and retention controls. AUC shows mean $\pm$ sample standard deviation; final scores are three-seed means.}\label{tab:controls-OBJECT}
\begin{tabular}{lrrr}
\toprule
Method & AUC $\uparrow$ & OldFinal $\uparrow$ & FinalAll $\uparrow$ \\
\midrule
Expert control $C$ & $85.00\pm3.07$ & 77.22 & 82.50 \\
Optimizer-matched $C$ & $80.90\pm1.20$ & 72.78 & 77.50 \\
Ungated $Q$ & $84.62\pm2.49$ & 78.33 & 82.92 \\
Complete method $S$ & $86.49\pm2.19$ & 81.11 & 85.00 \\
\bottomrule
\end{tabular}
\end{table}

Against optimizer-call-matched $C$, the complete method improves mean AUC
by 3.54 and 5.59 points on GOAL and OBJECT. Increasing $C$ updates alone
therefore does not explain the full advantage. Against ungated $Q$, mean
AUC increases by 3.75 and 1.88 points; their physical optimizer budgets differ.
\subsection{Sensitivity, behavior, and resources}
\begin{table}[tb]
\centering\small
\caption{GOAL retention-weight sensitivity: three-seed means in percent. The primary 0.10 setting reuses accepted historical runs; additional settings use the same scientific budgets, with hardware differences recorded.}\label{tab:lambda}
\begin{tabular}{rrrrrr}
\toprule
$\lambda_R$ & AUC & FWT & NBT $\downarrow$ & OldFinal & FinalAll \\
\midrule
0.05 & 83.02 & 85.42 & 4.63 & 79.44 & 81.67 \\
0.10 & 87.26 & 89.58 & 5.28 & 83.89 & 85.00 \\
0.20 & 83.23 & 82.92 & -0.93 & 85.00 & 82.92 \\
\bottomrule
\end{tabular}
\end{table}

The primary weight 0.1 achieves the highest mean AUC and FinalAll. A weight of 0.2 increases OldFinal but reduces FWT, revealing a retention--acquisition trade-off. These results do not establish invariance to the retention weight.

\begin{figure}[tb]\centering
\fbox{\parbox[c][25mm][c]{0.88\linewidth}{\centering
Figure slot: stage-wise success curves from accepted matrices\\
Show both suites and seed variability}}
\caption{Planned learning-curve figure. No synthetic experimental curve
is included in this framework.}\end{figure}
\draftnote{Add paired SF/FS, task matrices, component diagnostics, and the
relationship between offline and closed-loop results.}
Across both suites and three seeds, the complete method uses 544800 accepted
updates and 1012500 production optimizer calls. Matched $C$ uses 1012500
calls; $Q$ uses 544800. Both $Q$ and the complete method use 240 training
collection episodes across six runs; this $C$ control uses none.
Budget matching concerns optimizer calls, not identical interaction or wall time.
\draftnote{Add storage, probe-forward, candidate-update, and wall-clock costs.}
